\documentclass[11pt]{article}

\usepackage[margin=1in]{geometry}
\usepackage{amsmath,amssymb,amsthm}
\usepackage{mathpartir}
\usepackage{hyperref}

\newtheorem{theorem}{Theorem}

\newcommand{\ty}[3]{#1 \vdash #2 : #3}
\newcommand{\cv}[4]{#1 \vdash #2 = #3 : #4}
\newcommand{\U}{\mathsf{U}}
\newcommand{\Id}{\mathsf{Id}}
\newcommand{\Rf}{\mathsf{Ref}}
\newcommand{\Jelim}{\mathsf{J}}
\newcommand{\redone}{\rightsquigarrow_1}
\newcommand{\red}{\rightsquigarrow^{*}}
\newcommand{\Ga}{\Gamma}
% domain (finite-element) codes, in bold roman, distinct from the syntax
\newcommand{\bU}{\mathbb{U}}
\newcommand{\cId}{\mathbf{Id}}
\newcommand{\rEl}{\mathbf{r}}
\newcommand{\sem}[1]{[\![#1]\!]}
\newcommand{\ev}[2]{#1 \lhd \sem{#2}_{\rho}}   % finite element approximates value of a term
\newcommand{\mem}[2]{#1 \in #2}                 % finite-element membership
\newcommand{\valid}[4]{#1 \Vdash #2 : #3 \;[\,#4\,]}

\title{The identity type $\Id\,A\,M\,N$ with $\Rf$ and $\Jelim$\\[2pt]
  \large typing, reduction, injectivity, and subject reduction}
\author{Agda development \texttt{ID/} (domain-semantics)}
\date{}

\begin{document}
\maketitle

\noindent
This note records the rules for the element-level Martin-L\"of identity type
added to the type-in-type theory ($\U:\U$) of the \texttt{ID/} development, and
states the two metatheorems proved for it. Everything below is machine-checked
in Agda (\texttt{-{}-without-K -{}-exact-split}) with no postulates, no holes,
and no termination/positivity pragmas.

\medskip
\noindent
The eliminator is Martin-L\"of's \emph{original} $\Jelim$ with a
fully dependent motive. Writing the motive and base types out,
\[
  C : (X\,Y:A)\to \Id\,A\,X\,Y \to \U,
  \qquad
  D : (X:A)\to C\,X\,X\,(\Rf\,X),
\]
so that $\Jelim\,C\,D\,P$ at $P:\Id\,A\,M\,N$ has type $C\,M\,N\,P$ (the
application spine $((C\,M)\,N)\,P$).

\paragraph{Notation.}
Syntactic terms and types are written with the capital and standard names of the
rules below — a base type $A$, endpoint terms $M,N$, a motive $C$, a base term
$D$, a proof term $P$, bound variables $X,Y$. Finite elements of the semantic
domain are written with the \emph{reserved lowercase letters} $t,u,v,w,c$
(together with $\bot$, the universe code $\bU$, and the codes $\cId(t,u,v)$ and
$\rEl(w)$). We write $\sem{L}_\rho$ for the value of a term $L$ in an environment
$\rho$, and $\ev{u}{L}$ for ``$u$ is a finite approximant of that value''.

\section*{Typing rules}

\begin{mathpar}
\inferrule[Id]{\ty{\Ga}{A}{\U} \\ \ty{\Ga}{M}{A} \\ \ty{\Ga}{N}{A}}
  {\ty{\Ga}{\Id\,A\,M\,N}{\U}}
\and
\inferrule[Ref]{\ty{\Ga}{A}{\U} \\ \ty{\Ga}{M}{A}}
  {\ty{\Ga}{\Rf\,M}{\Id\,A\,M\,M}}
\end{mathpar}

\begin{mathpar}
\inferrule[J]{
  \ty{\Ga}{A}{\U} \\ \ty{\Ga}{M}{A} \\ \ty{\Ga}{N}{A} \\\\
  \ty{\Ga}{C}{(X\,Y:A)\to \Id\,A\,X\,Y\to\U} \\
  \ty{\Ga}{D}{(X:A)\to C\,X\,X\,(\Rf\,X)} \\
  \ty{\Ga}{P}{\Id\,A\,M\,N}}
  {\ty{\Ga}{\Jelim\,C\,D\,P}{C\,M\,N\,P}}
\end{mathpar}

\section*{Conversion rules}

\begin{mathpar}
\inferrule[J-$\beta$]{
  \ty{\Ga}{A}{\U} \\ \ty{\Ga}{M}{A} \\
  \ty{\Ga}{C}{\dots} \\ \ty{\Ga}{D}{\dots}}
  {\cv{\Ga}{\Jelim\,C\,D\,(\Rf\,M)}{D\,M}{C\,M\,M\,(\Rf\,M)}}
\end{mathpar}

\noindent
$\Jelim$ fires only on a \emph{literal} $\Rf\,M$, whose type is the diagonal
$\Id\,A\,M\,M$; both sides of J-$\beta$ therefore have the \emph{same} type
$C\,M\,M\,(\Rf\,M)$, so the rule carries no endpoint or motive equalities
and needs no injectivity. The congruences are

\begin{mathpar}
\inferrule[Id-cong]{
  \cv{\Ga}{A}{A'}{\U} \\ \cv{\Ga}{M}{M'}{A} \\ \cv{\Ga}{N}{N'}{A}}
  {\cv{\Ga}{\Id\,A\,M\,N}{\Id\,A'\,M'\,N'}{\U}}
\and
\inferrule[Ref-cong]{\cv{\Ga}{M}{M'}{A}}
  {\cv{\Ga}{\Rf\,M}{\Rf\,M'}{\Id\,A\,M\,M}}
\end{mathpar}

\begin{mathpar}
\inferrule[J-cong]{
  \cv{\Ga}{C}{C'}{\dots} \\ \cv{\Ga}{D}{D'}{\dots} \\ \cv{\Ga}{P}{P'}{\Id\,A\,M\,N}}
  {\cv{\Ga}{\Jelim\,C\,D\,P}{\Jelim\,C'\,D'\,P'}{C\,M\,N\,P}}
\end{mathpar}

\noindent
(Grey premises re-typing the components are elided; see
\texttt{ID/Syntax/Typing.agda}.)

\section*{Head reduction}

Reduction is untyped weak head reduction. The single-step relation
$L \redone L'$ has the following $\Id$-clauses (alongside the usual
$\beta$ and application-congruence for functions):

\begin{mathpar}
\inferrule[J-red]{ }{\Jelim\,C\,D\,(\Rf\,M) \;\redone\; D\,M}
\and
\inferrule[J-scrut]{P \redone P'}{\Jelim\,C\,D\,P \;\redone\; \Jelim\,C\,D\,P'}
\end{mathpar}

\noindent
\textsc{J-red} is the $\iota$-step; \textsc{J-scrut} lets reduction proceed into
the proof position so a non-canonical proof can reduce toward a $\Rf$. Each step
is a typed conversion (\textsc{J-red} $\mapsto$ J-$\beta$, \textsc{J-scrut}
$\mapsto$ J-cong), and $\redone$ is deterministic.

\section*{Domain semantics of $\Id$ and $\Jelim$}

The model is the poset of \emph{finite elements} $t,u,v,w,c,\dots$ of a Scott
domain, with an information order $u\le v$, a bottom $\bot$, a universe code
$\bU$, function elements, and $\Pi$-codes. The identity type adds two shapes:
\[
  \cId(t,u,v) \quad\text{(the \emph{code} of a type $\Id\,A\,M\,N$: $t$ codes
  $A$, $u$ codes $M$, $v$ codes $N$),}
\]
\[
  \rEl(w) \quad\text{(a \emph{proof value}: a witness $w$, the model's analogue
  of $\lambda$ for functions).}
\]
Semantic membership is written $\mem{u}{a}$ (``the finite element $u$ is a
member of the type-code $a$''). The characteristic clause for the identity type
is
\[
  \mem{\rEl(w)}{\cId(t,u,v)}
  \iff
  \mem{w}{t}\ \wedge\ w\le u\ \wedge\ w\le v ,
\]
i.e.\ a proof of $\cId(t,u,v)$ is exactly a witness of the underlying type that
lies \emph{below both endpoints} — the domain-theoretic reading of ``$u$ and
$v$ agree down to $w$''. All mismatched shapes (e.g.\ $\mem{\rEl(w)}{\Pi(\dots)}$
or $\mem{\cId(\dots)}{\Pi(\dots)}$) are empty, so identity codes and proof
values interact only with one another.

\medskip
\noindent
The approximation relation $\ev{u}{L}$ (``the finite element $u$ approximates
$\sem{L}_\rho$'') has these clauses for the new formers (for $u\neq\bot$; every
term is approximated by $\bot$):
\begin{align*}
  \ev{\cId(t,u,v)}{\Id\,A\,M\,N}
    &\iff \ev{t}{A}\ \wedge\ \ev{u}{M}\ \wedge\ \ev{v}{N}, \\[2pt]
  \ev{\rEl(w)}{\Rf\,M}
    &\iff \ev{w}{M}, \\[2pt]
  \ev{c}{\Jelim\,C\,D\,P}
    &\iff \exists\,w.\ \ev{w}{P}\ \wedge\ \mathsf{br}(c,w),
\end{align*}
where the branch condition on a proof value is
$\;\mathsf{br}(c,\rEl(w')) \iff \ev{\,\mathsf{fun}(w'\mapsto c)\,}{D}\;$ (and is
empty when the proof value is not a $\Rf$).
So the value of $\Id\,A\,M\,N$ is assembled componentwise into an $\cId$-code,
that of $\Rf\,M$ into a proof value, and the value of $\Jelim\,C\,D\,P$
dispatches on the proof: whenever a proof value $\rEl(w')$ approximates
$\sem{P}_\rho$, the value $\sem{D}_\rho$ applied to the witness $w'$ approximates
$\sem{\Jelim\,C\,D\,P}_\rho$ (the single-edge condition ``$D$ sends $w'$ to
$c$''). \textbf{The motive $C$ does not occur in the value} — it constrains only
the \emph{type}. This is the based-J computation rule read in the model.

\section*{The logical relation and adequacy}

Adequacy is proved through a Kripke logical relation between the syntax and the
model. For each type it defines, at every finite element, when a term is
\emph{valid}; write $\valid{\Ga}{L}{A}{\mem{u}{a}}$ for ``$L$ is valid at type
$A$ and code $a$, with value $u$''. At an identity type the relation unfolds to:
\[
\begin{aligned}
  &\valid{\Ga}{P}{\Id\,A\,M\,N}{\mem{\rEl(w)}{\cId(t,u,v)}} \iff{} \\
  &\qquad P \red \Rf\,P_0 \ \text{for some term $P_0$ with }
  \cv{\Ga}{P_0}{M}{A}\ \text{and}\ \cv{\Ga}{P_0}{N}{A}\ \text{(both valid).}
\end{aligned}
\]
That is: a valid proof \emph{head-reduces to a canonical $\Rf\,P_0$} whose
witness term $P_0$ is validly equal to both endpoints. (Because proofs are
proof-irrelevant, the underlying $\cId$-code carries no selection edge; the two
endpoint equalities are the whole content.)

\medskip
\noindent
\textbf{Adequacy.} By induction on derivations:
\[
  \ty{\Ga}{L}{A} \ \Longrightarrow\ L \text{ valid}, \qquad
  \cv{\Ga}{L}{L'}{A} \ \Longrightarrow\ L,L' \text{ validly equal.}
\]

\noindent
\emph{The $\Jelim$ case} (the based-J driver). Suppose
$\ty{\Ga}{\Jelim\,C\,D\,P}{C\,M\,N\,P}$ and let $c$ be a finite element
approximating $\sem{\Jelim\,C\,D\,P}_\rho$.

\begin{enumerate}\itemsep2pt
\item By the approximation clause, $c$ is accompanied by a semantic proof value:
  a $\rEl(w)$ approximates $\sem{P}_\rho$, and $\mem{\rEl(w)}{\cId(t,u,v)}$, so
  $w\le u$ and $w\le v$.
\item The induction hypothesis for the proof $P$ at $\Id\,A\,M\,N$ gives, via the
  relation above, a head reduction $P \red \Rf\,P_0$ whose witness term $P_0$ is
  validly equal to $M$ and to $N$. Hence
  \[
    \Jelim\,C\,D\,P \ \red\ \Jelim\,C\,D\,(\Rf\,P_0) \ \redone\ D\,P_0
    \qquad\text{(J-scrut${}^*$ then J-$\beta$).}
  \]
\item \emph{$D$-edge.} The induction hypothesis for the base $D$ (a valid
  function), applied to the witness, makes $D\,P_0$ valid at the diagonal type
  $C\,P_0\,P_0\,(\Rf\,P_0)$.
\item \emph{Type transport.} From $\cv{\Ga}{P_0}{M}{A}$, $\cv{\Ga}{P_0}{N}{A}$
  and $\Rf\,P_0 = P$ the syntactic types are convertible,
  $\cv{\Ga}{C\,P_0\,P_0\,(\Rf\,P_0)}{C\,M\,N\,P}{\U}$; together with monotonicity
  of the type-value along the information order this carries the validity of
  $D\,P_0$ to the goal type $C\,M\,N\,P$.
\item \emph{Head expansion.} Finally, validity is stable under the head reduction
  $\Jelim\,C\,D\,P \red D\,P_0$ of step~2 (the contractum is valid, and the redex
  converts to it by J-$\beta$), so $\Jelim\,C\,D\,P$ is valid at $C\,M\,N\,P$.
  \hfill$\square$
\end{enumerate}

\noindent
The two-substitution (equality) version of this argument handles J-cong; it
splits the two witnesses coming from the two sides and bridges them with the
endpoint equalities.

\section*{Metatheorems (corollaries of adequacy)}

Both results are obtained by running adequacy at the \emph{bottom environment}
and reading the logical relation off the model.

\begin{theorem}[Id-injectivity, \texttt{ID/IdInjectivity.agda}]
If $\;\cv{\Ga}{\Id\,A_0\,M_0\,N_0}{\Id\,A_1\,M_1\,N_1}{\U}\;$ then
\[
  \cv{\Ga}{A_0}{A_1}{\U},\qquad
  \cv{\Ga}{M_0}{M_1}{A_0},\qquad
  \cv{\Ga}{N_0}{N_1}{A_0}.
\]
\end{theorem}

\noindent
Evaluate the given conversion at the bottom environment: it yields a validated
equality at the type-code $\cId(\bot,\bot,\bot)$, whose stored data are exactly
the component conversions between domain, left and right endpoints; a
head-reduction uniqueness lemma pins the codes, giving the three conversions.

\begin{theorem}[Subject reduction, \texttt{ID/SubjectReduction.agda}]
If $\;\ty{\Ga}{L}{A}\;$ and $\;L \redone L'\;$ then $\;\ty{\Ga}{L'}{A}$.
\end{theorem}

\noindent
For the \textsc{J-red} step the redex is $\Jelim\,C\,D\,(\Rf\,K)$ with
$\Rf\,K:\Id\,A\,M\,N$; the contractum $D\,K$ already has type
$C\,K\,K\,(\Rf\,K)$, retyped to the goal $C\,M\,N\,(\Rf\,K)$ using $K=M$ and
$K=N$ extracted from $\Rf\,K:\Id\,A\,M\,N$. For \textsc{J-scrut} the result type
$C\,M\,N\,P$ is transported along $P=P'$, obtained from reduction-is-conversion;
the $\Id$-typed proof needs only the diagonal J-$\beta$, never
$\Id$-injectivity. The $\beta$-step for functions inverts $\lambda$-typing
through $\Pi$-injectivity, as in the base theory.

\medskip
\noindent
The adequacy theorems are in \texttt{ID/Adequacy/Value.agda}
(\texttt{adequacySub2} / \texttt{adequacyEqSub2}); the $\Jelim$ driver of steps
1--5 above is \texttt{ID/Adequacy/JApp.agda} (value), \texttt{JAppCross.agda}
and \texttt{JAppCongr.agda} (equality).

\end{document}
